\UnnumberedChapter{LỜI MỞ ĐẦU}

Trong kỷ nguyên bùng nổ của công nghệ viễn thông số và mạng di động toàn cầu, dịch vụ tin nhắn ngắn (Short Message Service --- SMS) vẫn tiếp tục giữ vai trò là một trong những phương thức giao tiếp trực tiếp, tức thời và có độ bao phủ người dùng cao nhất trên thế giới. Ngay cả khi các ứng dụng nhắn tin qua giao thức Internet (Over-The-Top --- OTT) như Zalo, Messenger, Telegram hay WhatsApp phát triển vượt bậc, SMS vẫn là kênh xương sống không thể thay thế đối với các dịch vụ xác thực bảo mật đa yếu tố (MFA/OTP), thông báo tài chính ngân hàng, cảnh báo y tế khẩn cấp và giao dịch thương mại điện tử.

Tuy nhiên, song hành cùng sự tiện lợi và tin cậy đó, vấn nạn tin nhắn rác (SMS Spam) đã và đang trở thành một thách thức nhức nhối đối với cả người dùng cá nhân lẫn các nhà mạng viễn thông. Không dừng lại ở việc phát tán các nội dung quảng cáo dịch vụ trái phép, tiếp thị làm phiền và gây tiêu tốn tài nguyên hạ tầng băng thông, tin nhắn rác hiện nay ngày càng bị các tổ chức tội phạm mạng lợi dụng để thực hiện các chiến dịch lừa đảo tài chính tinh vi, mạo danh cơ quan nhà nước, phát tán liên kết độc hại nhằm đánh cắp thông tin thẻ tín dụng, chiếm đoạt tài khoản định danh số và mã độc gián điệp.

Đứng trước tính chất nghiêm trọng của vấn đề, các phương pháp lọc tin nhắn truyền thống dựa trên luật tĩnh (Rule-based Filtering) hoặc lập danh sách đen các số thuê bao (Blacklist Filtering) đã nhanh chóng bộc lộ sự bất lực. Những kẻ phát tán tin rác dễ dàng vượt qua các bộ lọc cứng bằng cách liên tục thay đổi số điện thoại gửi đi thông qua SIM rác, cố tình chèn ký tự lạ hoặc biến tấu chính tả văn bản. Do đó, việc nghiên cứu và ứng dụng các kỹ thuật Học máy (Machine Learning) và Xử lý ngôn ngữ tự nhiên (NLP) nhằm xây dựng hệ thống nhận diện và phân loại tin nhắn rác tự động, có khả năng học hỏi các quy luật tiềm ẩn từ dữ liệu thực tế và thích ứng linh hoạt với các mẫu tin nhắn mới là một yêu cầu mang tính cấp thiết và có giá trị ứng dụng thực tiễn to lớn.

Trong khuôn khổ học phần \textbf{Nhập môn Học máy (mã học phần: FIT4013)} tại Khoa Công nghệ Thông tin --- Trường Đại học Đại Nam, dưới sự chỉ dẫn khoa học của ThS. Lê Thị Thùy Trang, Nhóm 13 đã lựa chọn triển khai \textbf{Đề tài 3: Phân loại tin nhắn rác (SMS Spam Collection)}. Mục tiêu cốt lõi của đề tài là vận dụng một cách hệ thống và bài bản toàn bộ các kiến thức từ tiền xử lý dữ liệu thô, trích xuất đặc trưng văn bản, thiết kế pipeline kiểm soát rò rỉ dữ liệu (Data Leakage), đến việc huấn luyện, tối ưu hóa siêu tham số và so sánh hiệu năng giữa ba thuật toán học máy tiêu biểu: \textbf{Multinomial Naive Bayes (MNB)}, \textbf{K-Láng giềng gần nhất (K-Nearest Neighbors --- K-NN)} và \textbf{Cây quyết định (Decision Tree --- DT)} trên bộ dữ liệu chuẩn quốc tế SMS Spam Collection.

Báo cáo bài tập lớn này được biên soạn công phu nhằm phản ánh toàn diện toàn bộ hành trình nghiên cứu thực nghiệm của nhóm, bao gồm 6 chương chính cùng các phần bổ trợ:
\begin{itemize}
    \item \textbf{Chương 1: Tổng quan về bài toán và cơ sở dữ liệu} --- Trình bày bối cảnh thực tiễn, mục tiêu nghiên cứu, nguồn gốc bộ dữ liệu SMS Spam Collection, quá trình khám phá dữ liệu chuyên sâu (EDA), phân tích mất cân bằng lớp và thiết lập hệ giả thuyết nghiên cứu.
    \item \textbf{Chương 2: Cơ sở lý thuyết và các thuật toán học máy áp dụng} --- Hệ thống hóa nền tảng toán học và thuật toán của Xử lý ngôn ngữ tự nhiên (BoW, TF-IDF, n-gram), nguyên lý toán xác suất của Multinomial Naive Bayes, cơ chế khoảng cách của K-NN, giải thuật chia nhánh của Cây quyết định và các độ đo đánh giá phân loại.
    \item \textbf{Chương 3: Thiết kế giải pháp và quy trình học máy (Pipeline)} --- Chi tiết hóa kiến trúc hệ thống, chiến lược phân tách tập dữ liệu chống rò rỉ thông tin, 3 luồng tiền xử lý (P0, P1, P2), 4 không gian biểu diễn đặc trưng, cơ chế tối ưu hóa siêu tham số GridSearchCV 5-fold CV và cấu trúc module mã nguồn.
    \item \textbf{Chương 4: Thực nghiệm và đánh giá kết quả} --- Trình bày chi tiết các kịch bản thực nghiệm từ baseline đến mô hình tinh chỉnh, khảo sát độ nhạy của hệ số làm trơn $\alpha$, độ nhạy tham số vectorizer, đối sánh toàn diện các chỉ số (Accuracy, Precision, Recall, F1, ROC-AUC) trên tập kiểm thử độc lập và kiểm tra độ ổn định qua nhiều seed.
    \item \textbf{Chương 5: Phân tích lỗi và thảo luận kỹ thuật} --- Mổ xẻ chi tiết từng ca dự đoán sai (Dương tính giả FP và Âm tính giả FN), phân tích mẫu lỗi chung của cả 3 mô hình, định lượng độ quan trọng của từ vựng qua log-odds ratio và feature importance, kiểm chứng hệ giả thuyết và chỉ ra các cạm bẫy dữ liệu (Data Artefacts).
    \item \textbf{Chương 6: Xây dựng ứng dụng demo và hướng phát triển} --- Giới thiệu ứng dụng web tương tác thời gian thực xây dựng bằng Streamlit với tính năng giải thích quyết định của mô hình (Explainable AI), bộ 8 kịch bản kiểm thử mẫu và định hướng phát triển hệ thống trong tương lai.
    \item \textbf{Phần cuối và Phụ lục} --- Kết luận tổng quan, Danh mục tài liệu tham khảo, Khai báo sử dụng AI minh bạch và Hướng dẫn cài đặt/vận hành chi tiết.
\end{itemize}

Báo cáo hy vọng sẽ là một tài liệu kỹ thuật hoàn chỉnh, đáp ứng xuất sắc các tiêu chí đánh giá khắt khe trong đề cương môn học và là minh chứng rõ nét cho năng lực làm chủ công nghệ của sinh viên.